\chapter{Work--Energy Equation}

\section*{Introduction}

Car brakes convert kinetic energy into heat through friction---the work done by friction equals the kinetic energy the car loses. A hydroelectric dam converts gravitational potential energy into kinetic energy of water, then into electrical energy through a turbine. The work-energy equation tracks these energy transformations: how much energy goes in, how much comes out, and how much is lost to friction.

\begin{enumerate}
\item \textbf{Work--Energy Equation:} The net work done equals the change in kinetic energy.

\item \textbf{Conservative Force and Potential Energy:} For conservative forces, work equals the negative change in potential energy.

\end{enumerate}
\section*{Fundamental Equation Used}

\[
W_{\text{net}} = \Delta K = \int_{\mathbf{r}_1}^{\mathbf{r}_2} \mathbf{F}_{\text{net}} \cdot d\mathbf{r}
\]


\section*{Assumptions}

\textbf{Assumptions:} The force $\mathbf{F}_{\text{net}}$ is the total (net) force. The particle moves along a definite path from $\mathbf{r}_1$ to $\mathbf{r}_2$. The mass is constant.


\textbf{Limitations:} Does not directly account for time-dependent forces (use impulse-momentum instead). For variable mass (rockets), the general form $\mathbf{F} = d\mathbf{p}/dt$ must be used. Does not include relativistic effects ($K = (\gamma - 1)mc^2$ at high speeds).


\section*{Mathematical Derivation}

\textbf{Step 1: Newton's Second Law as Starting Point.}

Consider a particle of mass $m$ moving along a path under a net force $\mathbf{F}_{\text{net}}$. Newton's second law states:

\begin{equation}
\mathbf{F}_{\text{net}} = m\,\frac{d\mathbf{v}}{dt} = m\,\mathbf{a}.
\end{equation}

We want to relate the force acting on the particle to the change in its kinetic energy.

\textbf{Step 2: Taking the Dot Product with Velocity.}

Dot both sides of Newton's second law with the velocity $\mathbf{v}$:

\begin{equation}
\mathbf{F}_{\text{net}}\cdot\mathbf{v} = m\,\frac{d\mathbf{v}}{dt}\cdot\mathbf{v}.
\end{equation}

The left side is the instantaneous power delivered to the particle. The right side can be rewritten using the identity:

\begin{equation}
\frac{d\mathbf{v}}{dt}\cdot\mathbf{v} = \frac{1}{2}\,\frac{d}{dt}(\mathbf{v}\cdot\mathbf{v}) = \frac{1}{2}\,\frac{d}{dt}(v^2) = v\,\frac{dv}{dt},
\end{equation}

where we used the product rule: $\frac{d}{dt}(\mathbf{v}\cdot\mathbf{v}) = 2\,\mathbf{v}\cdot\frac{d\mathbf{v}}{dt}$. Therefore:

\begin{equation}
\mathbf{F}_{\text{net}}\cdot\mathbf{v} = \frac{d}{dt}\!\left(\frac{1}{2}mv^2\right).
\end{equation}

\textbf{Step 3: Introducing Work and Kinetic Energy.}

The left side is the power, $P = \mathbf{F}_{\text{net}}\cdot\mathbf{v}$. The right side is the rate of change of the kinetic energy $K = \frac{1}{2}mv^2$. So:

\begin{equation}
P = \frac{dK}{dt}.
\end{equation}

Now integrate both sides over time from $t_1$ to $t_2$:

\begin{equation}
\int_{t_1}^{t_2} P\,dt = \int_{t_1}^{t_2} \frac{dK}{dt}\,dt = K(t_2) - K(t_1).
\end{equation}

The left side is the total work done:

\begin{equation}
W_{\text{net}} = \int_{t_1}^{t_2} \mathbf{F}_{\text{net}}\cdot\mathbf{v}\,dt.
\end{equation}

\textbf{Step 4: Converting to a Path Integral.}

Since $\mathbf{v}\,dt = d\mathbf{r}$ (the displacement along the path), the work integral becomes:

\begin{equation}
W_{\text{net}} = \int_{\mathbf{r}_1}^{\mathbf{r}_2} \mathbf{F}_{\text{net}}\cdot d\mathbf{r}.
\end{equation}

Therefore:

\begin{equation}
\boxed{W_{\text{net}} = \Delta K = \frac{1}{2}mv_2^2 - \frac{1}{2}mv_1^2.}
\end{equation}

This is the \textbf{work--energy theorem}: the net work done on a particle equals its change in kinetic energy.

\textbf{Step 5: Conservative Forces and Potential Energy.}

If the net force is conservative (i.e., it can be written as $\mathbf{F} = -\nabla U$ for some potential energy function $U$), then:

\begin{equation}
W_{\text{net}} = \int_{\mathbf{r}_1}^{\mathbf{r}_2} (-\nabla U)\cdot d\mathbf{r} = -(U_2 - U_1) = -\Delta U.
\end{equation}

Substituting into the work-energy theorem:

\begin{equation}
-\Delta U = \Delta K \quad \Longrightarrow \quad K_1 + U_1 = K_2 + U_2.
\end{equation}

This is the \textbf{conservation of mechanical energy}:

\begin{equation}
\boxed{E = K + U = \text{constant} \quad \text{(for conservative forces)}.}
\end{equation}

\textbf{Step 6: Non-Conservative Forces.}

When non-conservative forces (friction, air resistance) are present, the total work includes both conservative and non-conservative parts:

\begin{equation}
W_{\text{nc}} + W_{\text{c}} = \Delta K.
\end{equation}

Since $W_{\text{c}} = -\Delta U$:

\begin{equation}
W_{\text{nc}} = \Delta K + \Delta U = \Delta E_{\text{mech}}.
\end{equation}

The work done by non-conservative forces equals the change in total mechanical energy. Friction ($W_{\text{nc}} < 0$) converts mechanical energy to thermal energy, reducing the mechanical energy of the system.

\section*{Characteristics of the Equation}

The work-energy equation is a scalar equation (no direction), making it often simpler than the vector form of Newton's second law. It holds in any reference frame (though the value of work depends on the frame). It is the integral form of $F = ma$---equivalent but sometimes more convenient.
\section*{Solution Methods}

Direct integration of $\mathbf{F} \cdot d\mathbf{r}$ along the path. Energy diagrams (plot $U(x)$ and find turning points). Conservation of energy for conservative systems. Power analysis: $P = \mathbf{F} \cdot \mathbf{v} = dW/dt$.
\section*{Verifications}

Work-energy theorem verified by measuring force and displacement (or velocity before and after). Pendulum experiments: convert potential to kinetic energy and back, verify conservation. Inclined plane experiments: measure work done by gravity and friction.
\section*{Applications}

Ballistics (projectile trajectories), vehicle dynamics (fuel efficiency, braking distance), structural engineering (energy absorption in crashes), sports science (optimizing throws and jumps), renewable energy (wind turbine power output).
\section*{What Richard Feynman Would Have Asked}

\subsection*{1. The Plain-English Translation}
\textit{``What is the work-energy equation really telling me?''}

The Core Meaning: Energy is accounting. The work-energy equation is a balance sheet: ``energy in'' minus ``energy out'' equals ``energy stored.'' If you push a box and do positive work, the box gains kinetic energy. If friction does negative work, the box loses kinetic energy. The total work (all sources) equals the change in kinetic energy.

The Metaphor: Think of kinetic energy as money in your bank account. Work is like deposits (positive work adds energy) and withdrawals (negative work removes energy). The work-energy equation says: total deposits minus total withdrawals equals the change in your balance. You don't need to track every individual transaction---just the net effect.

\subsection*{2. The Localized Mechanics}
\textit{``At a single point along the path, what is the work being done?''}

He would press you on the dot product. The instantaneous power is $P = \mathbf{F} \cdot \mathbf{v}$. If force and velocity are in the same direction, positive work is done (energy is added). If they're opposite, negative work is done (energy is removed). If they're perpendicular, no work is done---magnetic forces on charged particles, for example, do zero work because $\mathbf{F} \perp \mathbf{v}$.

The Smoothness Check: He would ask whether the work-energy equation is truly equivalent to Newton's second law. Yes---integrating $F = ma$ along the path gives $\int F \, dx = \Delta(\frac{1}{2}mv^2)$. They are the same physics in different mathematical packaging. But the energy version is often simpler because it's scalar.

\subsection*{3. Testing the Extreme Limits}
\textit{``What happens at extreme speeds or energies?''}

The Relativistic Limit: At speeds approaching $c$, $K = (\gamma - 1)mc^2$ instead of $\frac{1}{2}mv^2$. The work-energy equation still holds, but the kinetic energy formula changes. At $v = 0.87c$, $\gamma = 2$ and $K = mc^2$---the kinetic energy equals the rest mass energy.

The Quantum Limit: In quantum mechanics, the work-energy theorem becomes the expectation value of the Hamiltonian. The ``work'' is related to transitions between energy eigenstates. The uncertainty principle means that energy and time cannot both be precisely defined.

The Dissipative Limit: For pure friction (no conservative forces), all kinetic energy is converted to heat. The work done by friction equals the initial kinetic energy: $\frac{1}{2}mv^2 = \mu mgd$, giving stopping distance $d = v^2/(2\mu g)$. This is why doubling speed quadruples stopping distance.

\subsection*{4. Dimensionality and Geometry}
\textit{``How does the path affect the total work?''}

He would point out that for conservative forces (gravity, spring force), the path doesn't matter---only the endpoints. For non-conservative forces (friction, air resistance), the path matters: a longer, winding path through a friction field dissipates more energy. This is why racing cars take the shortest path (minimizing friction losses).

\subsection*{5. Cross-Disciplinary Connection}
\textit{``Where else does this work-energy principle appear?''}

The work-energy principle appears in: thermodynamics (first law: $\Delta U = Q - W$), economics (work = revenue, kinetic energy = profit), ecology (energy budgets of organisms), information theory (Landauer's principle: erasing one bit requires $kT \ln 2$ of work), and social science (``social work'' produces ``social energy''---change in societal state). The accounting principle---inputs minus outputs equals storage---is universal.
\section*{FAQ}

\textbf{Q: Can an object have kinetic energy but zero net work?}\\
\textbf{A:} Yes---if the net force is zero, kinetic energy is constant. The object moves at constant velocity. Work is being done by individual forces, but they cancel out (equal positive and negative work).

\textbf{Q: Why can't we use the work-energy equation for friction?}\\
\textbf{A:} We can---but friction is non-conservative, so the work depends on the path taken. A longer path means more work done by friction, even between the same two points. The work-energy equation still holds; it just doesn't simplify to conservation of energy.
\section*{Important Bibliography}

\begin{enumerate}
\item Goldstein, H., Poole, C., and Safko, J., \textit{Classical Mechanics}, 3rd ed. (Addison-Wesley, 2002), Chapters 1--2.
\item Feynman, R.P., Leighton, R.B., and Sands, M., \textit{The Feynman Lectures on Physics}, Vol. I (Basic Books, 2011), Chapters 4--8.
\item Coriolis, G.G., \textit{Th\'eorie relative du mouvement des machines} (Bachelier, 1835).
\item Hand, L.N. and Finch, J.D., \textit{Analytical Mechanics} (Cambridge University Press, 1998), Chapters 2--4.
\end{enumerate}


